% Generated by submissions/prepare.py from paper/main.tex; edit the source, not this file.


\section{Defect register of the studied pipeline (summary)}
\label{app:register}

A line-by-line review of the pipeline surfaced 15 harness defects (H-1--H-15: import-time
side effects, environment-dependent tool invocation, missing response logging, list-content
crashes on 2026 models, usage-logger stacking) and nine behavioral defects
(Table~\ref{tab:register}). All harness defects and behavioral defects B-1--B-8 are fixed in
the study instrument; B-10's underlying mechanism was later removed outright, so the failure
analyst no longer produces code at all and nothing the judge writes can ship. The
prompt templates were separately rewritten after review (Sec.~\ref{sec:prompts}). Identifiers are stable against the
repository's full register, which includes locations, failure scenarios, and dispositions;
that register never assigned the identifier B-9, so the nine behavioral defects carry the
identifiers B-1--B-8 and B-10, and Table~\ref{tab:register} has no B-9 row by construction.
% trace: src/vgap/pipeline/REVIEW.md (B-9 was never assigned; IDs kept stable)

\begin{table*}[t]
\centering
\caption{Behavioral defects found during review; disposition in the study instrument.}
\label{tab:register}
\small
\begin{tabularx}{\textwidth}{lXl}
\toprule
ID & Finding & Disposition \\
\midrule
B-1 & Template render errors silently become empty prompts upstream & fixed \\
B-2 & Hanging generated test aborts the whole problem run & fixed \\
B-3 & Judge pass/fail verdicts via substring scan of the entire response & fixed \\
B-4 & First-match score regex; ``Score: 85'' normalizes to a perfect 1.0 & fixed \\
B-5 & GREEN judge verdict threshold hardcoded at 0.7, ignoring configuration & fixed \\
B-6 & Two result classes with different validation semantics across judges & fixed \\
B-7 & Latent \texttt{NameError} in the RED retry fallback (unreachable) & fixed \\
B-8 & GREEN's execution gate tests one file; VERIFY tests the whole suite & fixed \\
B-10 & Judge-corrected code never ships; travels only as prompt text & mechanism removed \\
\bottomrule
\end{tabularx}
\end{table*}

\section{Prompt skeletons}
\label{app:prompts}

Abbreviated structure of the two RED-phase templates (the GREEN pair follows the same
pattern; full templates in the artifact). Anchoring blocks are marked in the source and
removed mechanically to create ablation variants.

\paragraph{RED generator (test author).}
\begin{small}
\begin{Verbatim}[breaklines,breakanywhere]
You are writing the test suite for a Python function BEFORE any
implementation exists (TDD, RED phase). Your tests are the only
specification the implementer will see [...]
## Task                     <- title, description, signatures
## Specification details    <- full problem statement
## Scenarios to cover       <- worked examples from the spec
[ANCHOR] ## MANDATORY: Canonical Examples from Specification
  For EVERY example: one test named test_canonical_example_N with
  the EXACT inputs and EXACT expected output + derivation comment.
## Rules  1. only specified behavior ... 5. plain pytest
## Reviewer feedback on your previous suite   <- retry only
## Output: one fenced ```python block, nothing else
\end{Verbatim}
\end{small}

\paragraph{RED judge (test reviewer).}
\begin{small}
\begin{Verbatim}[breaklines,breakanywhere]
Judge one thing only: does this suite faithfully verify the
specification? [...]
[ANCHOR] ## Check 0 - MANDATORY Canonical Example Coverage
  any missing/altered example => Final score: 0.00
## Checks  1. fidelity (recompute every expected value)
           2. no invented constraints  3. coverage  4. sanity
## Scoring: deduction schedule (0.30/0.30/0.20/0.10, floor 0)
## Output: findings; ADD:/REMOVE:/FIX: directives (apply-verbatim);
           last line exactly `Final score: <x.xx>`
\end{Verbatim}
\end{small}

\section{False-accept case studies}
\label{app:cases}

The following HumanEval examples share a pattern: the self-authored suite anchors itself to the
specification's worked examples (the anchoring block of App.~\ref{app:prompts} enforces one
\texttt{test\_canonical\_example} per example), the implementation passes these, and both miss
a behavior that the held-out oracle checks. Three HumanEval false accepts make this pattern
specific: two from the all-weak control and one from the cell with a frontier judge, each accepted
with a top GREEN-judge score. Code and tests are copied exactly from the committed run logs.

\paragraph{HumanEval/32 (\texttt{find\_zero}, control cell).} The spec requests a root of
a polynomial with an even number of coefficients (hence odd degree) under the stated
leading-coefficient assumption. The self-suite tests only the two docstring examples
(\texttt{find\_zero([1,2])} $\rightarrow -0.5$, \texttt{find\_zero([-6,11,-6,1])} $\rightarrow
1.0$). The implementation uses Newton's method starting from \texttt{x=0.0} for $100$
iterations. Newton-from-zero converges on both example polynomials, so the suite passes, but this
is not a general root-finder: when the iteration stalls (\texttt{dfx==0}, and it breaks) or
diverges, it returns a non-root. HumanEval\texttt{+}'s extra polynomials cause exactly
this, and the oracle rejects. The tests only confirm the two points the code was tailored for, not
the full specification.

\paragraph{HumanEval/74 (\texttt{total\_match}, control cell).} The spec returns the list with
the smaller total character count, \emph{and returns the first list if the totals are equal}.
The implementation handles the tie incorrectly: \texttt{if total1 < total2: return lst1} else
\texttt{return lst2}, so equal totals default to \texttt{lst2}. None of the five generated
example tests presents a clear tie---the only equal-total case is \texttt{([],[])}, where
both answers are \texttt{[]}, so the tie branch is never tested. HumanEval\texttt{+} includes
an equal-total pair with different contents; the code returns the second list where the spec
requires the first, and the oracle rejects.

\paragraph{HumanEval/153 (\texttt{Strongest\_Extension}, frontier-judge cell).} Strength is
calculated as (uppercase count) minus (lowercase count); the strongest extension is the one with
the highest value, with ties going to the earliest occurrence. The implementation sets the running
maximum to \texttt{0} at the start and only updates it when \texttt{strength > strongest\_strength}.
This means if every extension has a strength of zero or less, and the strongest is not the first
one, the loop never updates and it ends up returning \texttt{extensions[0]}. The two generated
tests do not trigger this scenario---example~1 is all-negative, but its first element is already
the strongest, so the bug is hidden. As a result, the suite passes, and the frontier judge gave it
a score of $1.0$. HumanEval\texttt{+} provides an all-negative case where the maximum is not the
first, and in that case, the oracle rejects it.

\paragraph{At application scale.} RGRBench also contains oracle-relative disagreements
whose source is the requirements contract. The
\texttt{strong\_judge} run on \texttt{numbers\_to\_words} certifies itself, but the held-out oracle
only passes $20$ out of its $22$ tests; App.~\ref{app:telephone} repeats this false accept
completely. The per-package artifacts (including self-authored tests, the generated package, and
the oracle's result) are committed for all $534$ runs.

\section{An RGRBench requirements--oracle mismatch: \texttt{numbers\_to\_words}}
\label{app:telephone}

The false accept of Figure~\ref{fig:misalignment}, laid out end to end. The pipeline sees only
the requirements below; from them it authors its tests and its implementation, and reports success
once the implementation passes those tests. The held-out oracle is applied afterward and never
enters the loop. The implementation passes the oracle's spelling tests; the observed disagreement is the
exception \emph{type} on the range check. The requirement fixes the message but not the
type, the self-test accepts any \texttt{Exception}, and the oracle demands
\texttt{ValueError}. This example supports an oracle-relative gap, but does not establish
that the agent misread or violated its supplied requirements. The historical input is
retained; silently adding the missing type now would change the evaluated task.

\paragraph{Requirements given to the pipeline.} User-story narratives are omitted; the acceptance
criteria are verbatim.
\begin{small}
\begin{Verbatim}[breaklines,breakanywhere]
# Spelling numbers as English words

## Overview
A number-spelling service turns any whole number from 0 through
9999 into its English words, following everyday conventions:
hyphenated compounds for twenty-one through ninety-nine, no "and"
between parts, and plain spaces between the thousands, hundreds,
and remainder portions.

### US-1: Spelling numbers up to twenty
- AC-1.1: Zero is spelled "zero", and single digits use their
  names, such as "five" and "eight".
- AC-1.2: Ten is spelled "ten", and the teens are single
  unhyphenated words, such as "thirteen" and "nineteen".

### US-2: Spelling two-digit numbers
- AC-2.1: Round tens are single words with no hyphen, such as
  "twenty" and "ninety".
- AC-2.2: Other two-digit numbers are hyphenated tens-units
  compounds: 21 is "twenty-one", 77 is "seventy-seven", and 99
  is "ninety-nine".

### US-3: Spelling hundreds
- AC-3.1: 100 is "one hundred"; the leading "one" is always kept.
- AC-3.2: A hundreds part is followed by the spelled remainder
  separated by a space: 303 is "three hundred three", 555 is
  "five hundred fifty-five", and 115 is "one hundred fifteen".
- AC-3.3: The word "and" never joins the hundreds to the
  remainder.

### US-4: Spelling thousands
- AC-4.1: 1000 is "one thousand" with the leading "one" kept,
  and round thousands follow the same pattern, such as "two
  thousand".
- AC-4.2: The thousands part is followed by the spelled
  remainder: 2400 is "two thousand four hundred", 3466 is "three
  thousand four hundred sixty-six", and an empty hundreds part is
  simply omitted, so 5005 is "five thousand five".
- AC-4.3: The largest supported number, 9999, is "nine thousand
  nine hundred ninety-nine".

### US-5: Rejecting numbers outside the range
- AC-5.1: Negative numbers are rejected with an error whose
  message includes "must be in 0..9999".
- AC-5.2: Numbers greater than 9999 are rejected with the same
  error message wording.
\end{Verbatim}
\end{small}

\paragraph{Tests the pipeline authored (range check).} It also wrote a passing test for every
spelling rule; the two range checks are the ones that decide this case. Each asks only for a
generic \texttt{Exception}.
\begin{small}
\begin{Verbatim}[breaklines,breakanywhere]
def test_reject_negative_numbers():
    with pytest.raises(Exception, match=r"must be in 0\.\.9999"):
        number_to_words(-1)                        # AC-5.1

def test_reject_numbers_greater_than_[...]ninety_nine():
    with pytest.raises(Exception, match=r"must be in 0\.\.9999"):
        number_to_words(10000)                     # AC-5.2
\end{Verbatim}
\end{small}

\paragraph{The generated implementation.} Digit tables and the remainder-spelling branch (both
correct) are omitted with \texttt{...} and \texttt{[...]}; all other lines are verbatim. The range
check raises a bare \texttt{Exception}.
\begin{small}
\begin{Verbatim}[breaklines,breakanywhere]
def number_to_words(n):
    if n < 0 or n > 9999:
        raise Exception("must be in 0..9999")
    if n == 0:
        return "zero"
    units = ["", "one", "two", ...]      # digit tables abbreviated
    teens = ["ten", "eleven", "twelve", ...]
    tens = ["", "", "twenty", "thirty", ...]
    thousands = n // 1000
    hundreds = (n % 1000) // 100
    remainder = n % 100
    words = []
    if thousands:
        words.append(units[thousands] + " thousand")
    if hundreds:
        words.append(units[hundreds] + " hundred")
    if remainder:
        [...]  # spell the 0..99 remainder, hyphenating 21..99
    return " ".join(words)
\end{Verbatim}
\end{small}

\paragraph{The held-out oracle (applied after the loop).} The other twenty tests cover the
spelling rules and all pass; the two range checks require \texttt{ValueError}, which a bare
\texttt{Exception} is not, so running the full oracle against the implementation gives $20$ of
$22$.
\begin{small}
\begin{Verbatim}[breaklines,breakanywhere]
def test_negative_numbers_are_rejected() -> None:
    with pytest.raises(ValueError, match="must be in 0..9999"):
        number_to_words(-1)

def test_numbers_above_9999_are_rejected() -> None:
    with pytest.raises(ValueError, match="must be in 0..9999"):
        number_to_words(10000)
\end{Verbatim}
\end{small}

